% Vietnamese translation for OI Public Library
% Source-PDF: IOI中国国家候选队论文/1999/石润婷.pdf
\documentclass[11pt]{article}
\usepackage{oipl-vn}

\title{Ẩn hóa, đa chiều hóa, mở hóa: bàn về tính linh hoạt của mô hình hóa toán học trong thi tin học hiện nay}
\author{Shi Runting}
\date{}

\begin{document}
\maketitle

\origtitle{Concealment, multidimensionalization, and openness: on the flexibility of mathematical modeling in today's informatics competitions}
\sourcepath{IOI China candidate team papers / 1999 / Shi Runting.pdf}

\paragraph{Từ khóa.}
Mô hình hóa toán học; ẩn hóa; đa chiều hóa; mở hóa.

\section*{Tóm tắt}

Mô hình hóa toán học là một bộ phận hữu cơ của Olympic Tin học. Các kỳ thi
tin học hiện nay ngày càng theo đuổi tính linh hoạt trong mô hình hóa toán
học. Biểu hiện của nó có thể khái quát thành ba xu hướng: mô hình được ẩn
đi, được mở rộng theo nhiều chiều, và được mở ra khỏi các khuôn mẫu cổ điển.
Bài viết này thông qua việc thảo luận ý nghĩa và biểu hiện của ba xu hướng
ấy để nghiên cứu các chiến lược giải bài tương ứng.

\section{Dẫn nhập}

Mô hình hóa toán học, với tư cách một bộ phận không thể thiếu của Olympic Tin
học, đã luôn tiến hóa, hoàn thiện và phát triển từ khi hình thức thi này ra
đời. Các kỳ thi tin học hiện nay ngày càng theo đuổi tính linh hoạt của mô
hình hóa toán học. Chính tính linh hoạt ấy làm vẻ hấp dẫn của mô hình hóa
toán học được thể hiện đầy đủ, đồng thời trao cho thi tin học sức sống lâu
dài và triển vọng phát triển rộng mở.

Qua việc quan sát và nghiên cứu một loạt bài thi mới, tôi nhận thấy tính linh
hoạt của mô hình hóa toán học trong thi tin học hiện nay có thể được khái
quát thành ba điểm: ẩn hóa, đa chiều hóa và mở hóa. Dưới đây, ta sẽ thông qua
việc bàn về ý nghĩa và biểu hiện của ba xu hướng ấy để rút ra các chiến lược
giải bài tương ứng.

\section{Nội dung chính}

\subsection{Ẩn hóa}

\subsubsection{Định nghĩa}

Nghĩa gốc của ``ẩn'' là mượn sự vật khác để che lấp. Khi đặt vào thi tin học,
``sự vật khác'' và đối tượng bị ``che lấp'' có chỉ định cụ thể. Nhìn ngay từ
chủ đề của chúng ta, đối tượng bị che lấp chính là mô hình toán học; còn
``sự vật khác'' ở đây là những tình huống thực tế phức tạp, nhiều lớp và dễ
gây nhầm lẫn.

Vì vậy, định nghĩa của việc ẩn hóa mô hình toán học trong thi tin học là:
dùng những tình huống thực tế phức tạp để che lấp mô hình toán học.

\subsubsection{Biểu hiện}

Một biểu hiện lớn của ẩn hóa trong thi tin học là ``mô hình cũ, bộ mặt mới'':
vẫn dùng một mô hình quen thuộc, nhưng tạo ra một khung cảnh hoàn toàn mới
để chứa nó, qua đó khoác cho mô hình vốn lộ rõ một bộ áo mới và ``che chắn''
nó. Vì thế, cùng một mô hình khi xuất hiện trong các bài thi khác nhau thường
có diện mạo rất khó lường. Hai bài ``Tuyến du lịch tối ưu'' (NOI 1997) và
``Tên lửa kiểu mới'' là những ví dụ điển hình; phát biểu bài toán được đặt ở
các phụ lục của bài gốc.

Bài thứ nhất mô tả một mạng đường phố gồm ``đường có bóng cây'' và ``phố du
lịch'', còn bài thứ hai mô tả một vấn đề ``nổ tên lửa''. Nếu chỉ nhìn bề mặt,
hai bài này dường như chẳng liên quan gì đến nhau. Tuy nhiên, đằng sau hai
hiện tượng bên ngoài hoàn toàn khác nhau lại ẩn chứa cùng một nguyên mẫu:
bài toán tìm đoạn con liên tiếp ``lớn nhất'' trong một dãy một chiều, trong
đó ``lớn nhất'' nghĩa là tổng phần tử lớn nhất. Cụ thể, với dãy
\[
  A=\{a_1,a_2,a_3,\ldots,a_N\},
\]
cần tìm
\[
  \operatorname{Maxsum}
    =\max\left\{\sum_{k=i}^{j} a_k \mid 1\le i\le j\le N\right\}.
\]

Điểm chung của hai bài này là bản thân đề bài không trực tiếp phơi bày mô
hình toán học, mà thông qua việc mô tả cụ thể một tình huống thực tế phức tạp
để yêu cầu thí sinh tự quy nạp và trừu tượng hóa ra mô hình toán học. Vì vậy,
chúng thể hiện đầy đủ đặc trưng ẩn hóa của mô hình toán học trong thi tin
học.

\subsubsection{Chiến lược: phương pháp vén màn sương}

Những tình huống thực tế phức tạp thường tạo cho người đọc cảm giác như nhìn
hoa trong sương. Nhưng chỉ cần ta có cái nhìn sắc bén, xuyên qua các hiện
tượng bề mặt rồi nắm được bản chất của mô hình, bài toán thường sẽ được giải
quyết dễ dàng.

Hãy phân tích bài ``Tên lửa kiểu mới'' để thấy cụ thể quá trình suy nghĩ
``vén màn sương'' và đào sâu bản chất bài toán. Trước hết, ta không nên bị
các khái niệm bề ngoài như ``bán kính che chắn'', ``bán kính tấn công'',
``điểm dân cư'' và ``lô cốt'' làm rối. Cần chú ý các sự thật sau.

\begin{enumerate}
  \item Có thể đổi điểm của mỗi điểm dân cư thành số đối của nó. Khi đó tổng
  lợi ích của vụ nổ bằng tổng điểm của tất cả điểm dân cư cộng tổng điểm của
  tất cả lô cốt. Nhờ vậy, ta có thể nhìn điểm dân cư và lô cốt theo cùng một
  cách.

  \item Một khi ``bán kính che chắn'' và ``bán kính tấn công'' đã được xác
  định, việc một công trình có bị phá hủy hay không chỉ phụ thuộc vào khoảng
  cách từ nó đến tâm nổ, mà không phụ thuộc vào vị trí cụ thể của nó trên mặt
  phẳng chiến trường. Do đó khi đọc dữ liệu, ta chỉ cần lưu khoảng cách từ
  mỗi điểm đến tâm nổ, và loại bỏ tọa độ \(x,y\) cụ thể.

  \item Có thể sắp xếp các điểm theo thứ tự từ gần đến xa tâm nổ, đồng thời
  gộp các điểm có cùng khoảng cách đến tâm thành một điểm đại diện, với điểm
  số bằng tổng điểm của chúng.
\end{enumerate}

Sau các bước này, dữ liệu có thể trở thành một bảng một chiều như
\[
  \begin{array}{ccccc}
    5 & 6 & -20 & 6 & 3
  \end{array}
\]
và nhiệm vụ của ta chính là tìm đoạn con có tổng lớn nhất trong bảng một
chiều đó, tức bản chất của mô hình. Bài gốc kèm chương trình tham khảo
\texttt{missile.pas}.

Tóm lại, khi cầm đề bài, ta không nên vội bắt tay lập trình ngay. Trước hết
cần bình tĩnh suy nghĩ và phân tích, lọc bỏ đúng những phần vô dụng, gây mê
hoặc trong các thông tin liên quan, rồi tinh chế ra phần then chốt. Nhờ vậy
ta mới nắm được ``mô hình cũ'' ẩn sau ``bộ mặt mới''. Tục ngữ có câu: mài dao
không làm chậm việc đốn củi. Chỉ sau khi suy nghĩ chu đáo và sâu sắc, ta mới
có khả năng xuyên qua hiện tượng để thấy bản chất vấn đề. Lúc ấy mới bắt tay
lập trình thì sẽ nắm chắc trong lòng bàn tay và đạt hiệu quả gấp bội.

\subsection{Đa chiều hóa}

Đa chiều hóa mô hình là một biểu hiện khác của tính linh hoạt trong mô hình
hóa toán học ở thi tin học. Nó có thể chia thành hai loại lớn: đa chiều hóa
``thực'' và đa chiều hóa ``ảo''.

\subsubsection{Đa chiều hóa thực}

\paragraph{Ý nghĩa.}
Đa chiều hóa ``thực'', đúng như tên gọi, là sự đa chiều hóa thực sự, có thể
``nhìn thấy, chạm thấy''. Đây là nội hàm của nó. Ngoại diên của nó là mô hình
được mở rộng từ đường sang mặt, và từ mặt sang không gian.

Hãy xem hai mô hình sau để cảm nhận cụ thể ý nghĩa của việc mở rộng ``từ
đường sang mặt rồi sang không gian''.

\begin{itemize}
  \item Mô hình 1: cho một số điểm trong mặt phẳng, tìm hình tròn nhỏ nhất
  phủ tất cả các điểm ấy.
  \item Mô hình 2: cho một số điểm trong không gian, tìm hình cầu nhỏ nhất
  phủ tất cả các điểm ấy.
\end{itemize}

Có thể thấy mô hình 2 là sản phẩm thu được bằng cách đa chiều hóa mô hình 1.
Trên thực tế, loại mô hình đa chiều hóa này đã nhiều lần xuất hiện trong đề
thi và bài luyện tập; mô hình 2 ở trên chỉ là một ví dụ nhỏ.

\paragraph{Chiến lược: phương pháp hạ chiều.}
Xu hướng đa chiều hóa ``thực'' làm tăng yếu tố không gian cần xét, từ đó làm
tăng độ khó của việc giải mô hình.

Khi giải các bài như vậy, ta thường cần lần theo mạch suy nghĩ của người ra
đề: cùng tiến dần từng bước, bắt đầu từ vấn đề có số chiều thấp; sau khi tìm
được cách giải phù hợp cho vấn đề thấp chiều, ta mở rộng và tổng quát hóa nó,
từ đó thu được cách giải cho vấn đề nhiều chiều tương ứng. Đây thực chất là
tư tưởng ``hạ chiều''. Ưu điểm của nó là trước hết biến phức tạp thành đơn
giản, giúp ta tìm được điểm bắt đầu cho việc phân tích. Sau khi đã nắm được
cách giải mô hình thấp chiều, việc quay từ đơn giản về khó hơn để thu được
cách giải bài toán nhiều chiều sẽ trở nên tự nhiên.

Tư tưởng hạ chiều có nhiều cách vận dụng khác nhau trong từng bài.

\paragraph{Phương pháp tương tự.}
Lấy bài ``Phủ sóng vệ tinh'' của NOI 1997 làm ví dụ. Trở ngại đầu tiên trong
quá trình phân tích bài này có lẽ là độ khó không gian. Nhưng nếu trước mắt
ta không phải là một tình huống ba chiều, mà chỉ là một tình huống hai chiều
đơn giản, có lẽ ta sẽ tự tin hơn rất nhiều. Vậy tại sao không thử giải mô
hình hai chiều trước, biết đâu từ đó lại nhận được gợi ý? Thực ra mô hình hai
chiều của bài này, tức bài toán tính tổng diện tích được phủ bởi nhiều hình
chữ nhật đồng phẳng có thể chồng lên nhau, đã xuất hiện trong bài ``Tính diện
tích hình'' ở IOI lần thứ năm, nên chúng ta khá quen thuộc. Thuật toán có thể
mô tả như sau:

\begin{enumerate}
  \item Tiền xử lý: xóa tất cả hình chữ nhật bị chứa hoàn toàn trong hình
  khác.
  \item Rời rạc hóa mặt phẳng.
  \item Thống kê tổng diện tích của mọi ô rời rạc bị phủ.
\end{enumerate}

Sau đó, bằng phép tương tự, ta suy ra thuận lợi cách giải cho mô hình ba
chiều tương ứng:

\begin{enumerate}
  \item Tiền xử lý: xóa tất cả hình hộp bị chứa hoàn toàn trong hình hộp
  khác.
  \item Rời rạc hóa không gian.
  \item Thống kê tổng thể tích của mọi khối ô rời rạc bị ``phủ''.
\end{enumerate}

Phương án giải mô hình ba chiều này gần giống phương án giải mô hình hai
chiều, chỉ khác là khi thống kê cần thêm một số tối ưu hóa vì lý do hiệu
suất. Như vậy, bằng cách vận dụng tư tưởng hạ chiều, ta dựa trên gợi ý từ lời
giải mô hình hai chiều để hoàn thành trọn vẹn lời giải mô hình ba chiều. Bài
gốc kèm chương trình tham khảo \texttt{cover.pas}.

Ví dụ trên cho thấy mô hình nhiều chiều được sinh ra từ mô hình thấp chiều,
nên khó tránh khỏi việc ``kế thừa'' một số đặc trưng của mô hình thấp chiều.
Điều đó giúp ta có khả năng dùng phép tương tự để áp dụng trực tiếp mẫu giải
của mô hình thấp chiều, qua đó ``đỡ đẻ'' cho vấn đề nhiều chiều phức tạp hơn.

\paragraph{Phương pháp hạ chiều được thực thi thật sự.}
Trong phương pháp tương tự ở trên, toàn bộ quá trình vận dụng tư tưởng
``hạ chiều'' thực ra diễn ra trong đầu ta, còn đối tượng thao tác thực sự của
chương trình vẫn là mô hình nhiều chiều. Do đó, ``hạ chiều'' chưa thật sự
được thể hiện trong chương trình.

Tuy nhiều bài đa chiều hóa trong thi có thể dùng phép tương tự để áp dụng mẫu
giải thấp chiều sẵn có, ta cũng nên thấy rằng cách này không phải lúc nào
cũng thuận lợi. Hãy xem bài ``Thám hiểm vũ trụ''.

Mô hình một chiều của bài này chính là bài toán đã nhắc ở trên: tìm dãy con
liên tiếp có tổng lớn nhất trong một dãy một chiều. Nhưng mẫu giải của mô
hình một chiều, tức phương pháp một vòng lặp, không thể áp dụng trực tiếp vào
không gian ba chiều.

Với bài này, cách dễ nghĩ đến hơn cả là vét cạn, nhưng hiệu suất của nó cực
kỳ thấp. Để xác định một hình hộp con cần tổng cộng sáu biến: tọa độ góc
trái-dưới-trước \((a_1,b_1,c_1)\) và tọa độ góc phải-trên-sau
\((a_2,b_2,c_2)\). Vì vậy cần sáu vòng lặp. Ngay cả khi không tính các vòng
lặp mới để tính tổng điểm của hình hộp nằm bên trong sáu vòng lặp ấy, độ phức
tạp thời gian của thuật toán đã đạt quy mô \(N^6\), với \(N\le 50\). Đây rõ
ràng là một con số không thể chấp nhận, nên thuật toán như vậy không nên
dùng.

Ở đây, ta có thể dùng một cách ``hạ chiều'' thực sự được đưa vào chương trình
để giải bài này.

Cụ thể, giả sử hình hộp đang tìm nằm giữa các lớp từ \(a_1\) đến \(a_2\) theo
chiều trên-xuống. Để tìm hình hộp có điểm cao nhất nằm giữa hai lớp
\(a_1\) và \(a_2\), tức để xác định tiếp \(b_1,b_2,c_1,c_2\), ta trước hết
nên nén các lớp \(a_1\ldots a_2\) từ ba chiều xuống hai chiều bằng cách cộng
các giá trị tương ứng theo chiều đó.

Lúc này bài toán đã được chuyển thành một mô hình hai chiều: tìm một hình chữ
nhật có tổng điểm lớn nhất trong một bảng hai chiều. Dễ thấy hơn, hình chữ
nhật ``lớn nhất'' này khi khôi phục về ba chiều chính là hình hộp ``lớn
nhất'' nằm giữa các lớp \(a_1\ldots a_2\).

Để giải mô hình hai chiều trên, ta dùng lại cùng phương pháp: trước hết xác
định \(b_1,b_2\), rồi nén các lớp \(b_1\ldots b_2\) từ hai chiều xuống một
chiều. Khi đó có thể dùng mẫu giải một chiều đã biết để hoàn tất lời giải.

Như vậy, khi cuối cùng giải mô hình một chiều, ta đã tận dụng một thuật toán
tốt sẵn có và tiết kiệm được một vòng lặp, hạ độ phức tạp thời gian xuống
\(N^5\). Quan trọng hơn, ta có thể dùng quá trình hạ chiều để giảm số lần tính
tổng lặp lại. Khi đã hoàn thành việc xét các lớp \(a_1\ldots a_2\), và bắt
đầu nén các lớp \(a_1\ldots a_2+1\), ta có thể dựa phép cộng dồn trên bảng
nén đã có của \(a_1\ldots a_2\): chỉ cần cộng thêm lớp \(a_2+1\) vào các vị
trí tương ứng trong bảng nén ấy. Như vậy, ta sử dụng hiệu quả thông tin đã
thu được và tránh được rất nhiều tính toán lặp. Quá trình nén từ hai chiều
xuống một chiều cũng tương tự. So với vét cạn, việc thống kê của thuật toán
này được phân tán giữa các tầng vòng lặp, chứ không lồng sâu vào vòng lặp
cuối cùng. Nhờ vậy, xu hướng tăng đột ngột của độ phức tạp thuật toán được
kiểm soát hiệu quả. Bài gốc kèm chương trình tham khảo \texttt{explore.pas}.

Trong quá trình giải bài này, chúng ta đưa quá trình ``hạ chiều'' vào thực
thể thao tác của chương trình. Đây chính là ``hạ chiều được thực thi thật
sự'' mà ta nói đến. Qua đối chiếu với vét cạn có thể thấy cách hạ chiều này
không chỉ đem lại mạch suy nghĩ rõ ràng cho việc phân tích, mà còn thường tạo
ra những hiệu quả bất ngờ và tinh tế.

Do bản thân các bài đa chiều hóa rất đa dạng, cách vận dụng tư tưởng hạ chiều
còn nhiều nữa. Điều then chốt là phải dựa vào đặc thù của từng bài để vận
dụng linh hoạt. Vì giới hạn khuôn khổ, ở đây chỉ giới thiệu hai cách vận dụng
khá thường gặp, mong có tác dụng gợi mở.

\subsubsection{Đa chiều hóa ảo}

\paragraph{Định nghĩa.}
Trong đa chiều hóa ``thực'' ở trên, ``chiều'' dùng theo nghĩa gốc: yếu tố cấu
thành không gian. Trong đa chiều hóa ``ảo'', ý nghĩa của ``chiều'' có thể mở
rộng thành ``yếu tố cấu thành mô hình toán học'' nói chung.

Theo đó, ý nghĩa của đa chiều hóa ``ảo'' là sự tăng lên của các yếu tố cấu
thành mô hình toán học.

\paragraph{Biểu hiện và chiến lược.}
Khác với đa chiều hóa ``thực'', đa chiều hóa ``ảo'' thể hiện trong mô hình
toán học và chương trình của ta bằng cách tăng tham số giai đoạn, tham số
trạng thái, v.v. Đây vừa là hình thức biểu hiện của đa chiều hóa ``ảo'', vừa
là chiến lược giải bài tương ứng.

Điển hình nhất trong số đó là sự đa chiều hóa của mô hình quy hoạch động.
Phương pháp gán nhãn trong lý thuyết đồ thị có thể xem như một tối ưu hóa của
quy hoạch động, nên ở đây cũng được xếp vào quy hoạch động.

Ta biết quy hoạch động cổ điển được tạo bởi ba lớp vòng lặp: giai đoạn, trạng
thái và quyết định. Nhưng quy hoạch động hiện nay thường không còn bị giới
hạn trong mẫu ba vòng lặp cũ kỹ ấy. Bằng cách tăng số biến giai đoạn và biến
trạng thái, mô hình được xây dựng trên nhiều lớp vòng lặp hơn. Các bài tiêu
biểu gồm ``Mua sắm cửa hàng'' ở IOI lần thứ bảy, ``Trò chơi xếp khối'' của
NOI 1997 và ``Trò chơi Roger'' trong kỳ tuyển chọn đội Trung Quốc dự IOI
1998; đề bài được đặt trong phụ lục của bài gốc.

Để làm rõ vấn đề, ta lấy ``Trò chơi Roger'' làm ví dụ cụ thể về biểu hiện của
đa chiều hóa ``ảo''. Trước khi xem bài này, hãy xét một mô hình quen thuộc để
so sánh.

Có một bảng hai chiều kích thước \(A\times B\). Mỗi ô trong bảng được gán một
giá trị nguyên là \(-1\) hoặc nằm trong khoảng \(0\ldots255\). Ta cần đi từ ô
\((x_1,y_1)\) đến ô \((x_2,y_2)\), trên đường không được qua ô \(-1\). Tổng
các số trên mọi ô đi qua được gọi là chi phí của đường đi. Hãy tìm chi phí
nhỏ nhất trong mọi đường đi có thể.

Ai cũng biết mô hình này là một mô hình quy hoạch động cổ điển. Phương trình
quy hoạch động có thể viết là
\[
\operatorname{Mincst}[F_x,F_y]=
\begin{cases}
0, & \operatorname{Map}[F_x,F_y]\ne -1,\\
\operatorname{MaxInt}, & \operatorname{Map}[F_x,F_y]=-1,
\end{cases}
\]
trong đó \((F_x,F_y)\) là tọa độ điểm xuất phát. Với \(1\le i\le A\),
\(1\le j\le B\),
\[
\operatorname{Mincst}[i,j]=
\begin{cases}
\operatorname{MaxInt}, & \operatorname{Map}[i,j]=-1,\\
\min\left\{
\begin{array}{l}
\operatorname{Mincst}[i-1,j]\quad(1\le i-1\le A),\\
\operatorname{Mincst}[i,j-1]\quad(1\le j-1\le B),\\
\operatorname{Mincst}[i+1,j]\quad(1\le i+1\le A),\\
\operatorname{Mincst}[i,j+1]\quad(1\le j+1\le B)
\end{array}
\right\}+\operatorname{Map}[i,j],
& \operatorname{Map}[i,j]\ne -1.
\end{cases}
\]

Bây giờ hãy xem bài ``Trò chơi Roger'' đã ``đa chiều hóa'' mô hình trên như
thế nào. Cũng là đi từ ô \((x_1,y_1)\) đến ô \((x_2,y_2)\), nhưng so với mô
hình đầu tiên, bài ``Roger'' đưa thêm vào hàng nghìn trạng thái khác nhau do
vị trí các số trên sáu mặt của bản thân Roger khác nhau. Vì vậy, để thể hiện
sự khác biệt giữa các trạng thái ấy, phương trình quy hoạch động mới nhất
thiết phải đưa thêm tham luận mới, tức một ``chiều'' mới:
\[
\operatorname{Mincst}[F_x,F_y,k]=
\begin{cases}
0, & \operatorname{Map}[F_x,F_y]\ne -1,\\
\operatorname{MaxInt}, & \operatorname{Map}[F_x,F_y]=-1,
\end{cases}
\qquad 1\le k\le 720,
\]
trong đó \((F_x,F_y)\) là tọa độ điểm xuất phát. Với \(1\le i\le N\),
\(1\le j\le M\), \(1\le k\le 720\),
\[
\operatorname{Mincst}[i,j,k]=
\begin{cases}
\operatorname{MaxInt}, & \operatorname{Map}[i,j]=-1,\\
\min\left\{
\begin{array}{l}
\operatorname{Mincst}[i-1,j,k_1]+\operatorname{dice}[k_1]\operatorname{Map}[i,j]\quad(1\le i-1\le A),\\
\operatorname{Mincst}[i,j-1,k_2]+\operatorname{dice}[k_2]\operatorname{Map}[i,j]\quad(1\le j-1\le B),\\
\operatorname{Mincst}[i+1,j,k_3]+\operatorname{dice}[k_3]\operatorname{Map}[i,j]\quad(1\le i+1\le A),\\
\operatorname{Mincst}[i,j+1,k_4]+\operatorname{dice}[k_4]\operatorname{Map}[i,j]\quad(1\le j+1\le B)
\end{array}
\right\},
& \operatorname{Map}[i,j]\ne -1.
\end{cases}
\]
Ở đây, \(k_d\) là trạng thái thu được sau khi Roger ở trạng thái \(k\) lăn
một lần theo hướng thứ \(d\); \(\operatorname{dice}[k_d]\operatorname{Map}[i,j]\)
là chi phí của bước hiện tại. Bài gốc kèm chương trình tham khảo
\texttt{roger.pas}.

\subsubsection{Tiểu kết}

Dù là đa chiều hóa ``thực'' hay ``ảo'', trọng tâm kiểm tra đều là phẩm chất
tư duy của thí sinh và năng lực mở rộng, di chuyển tri thức. Vì vậy trong
luyện tập hằng ngày, ta đặc biệt cần chú ý liên tưởng và so sánh, học cách
nhìn thi tin học bằng con mắt phát triển. Không nên thỏa mãn với việc luyện
một bài cụ thể; sau khi giải xong bài đó, cần chú ý mở rộng và đào sâu, chủ
động tiến hành đa chiều hóa thay vì bị động đối phó với đa chiều hóa. Chỉ như
vậy mới có thể tiến lên một tầng cao hơn.

\subsection{Mở hóa}

\subsubsection{Ý nghĩa}

Nghĩa gốc của ``mở'' là phá bỏ lệnh cấm, phong tỏa và ràng buộc. Trong thi
tin học, chúng ta mở rộng nó thành ``phá bỏ ràng buộc của các mô hình có sẵn,
cổ điển''.

Nếu các quá trình ẩn hóa và đa chiều hóa đã nói ở trên được xây dựng trên cơ
sở các mô hình hiện có, thì mở hóa là một quá trình đổi mới triệt để, vì nó
hoàn toàn phá vỡ các mẫu cũ của mô hình toán học trong thi tin học.

\subsubsection{Biểu hiện}

Những mô hình mở hóa liên tục xuất hiện trong thi đấu thường tạo cho người ta
cảm giác mới mẻ hoàn toàn. Đồng thời, chúng cũng thường làm các thuật toán
``có tên có tuổi'' không còn chỗ dùng. Vì mất đi sự tham chiếu từ mô hình
kinh điển, người giải thường khó tìm được điểm vào. Điều này kiểm tra đầy đủ
năng lực sáng tạo của thí sinh. Nó không chỉ yêu cầu ta biết linh hoạt vận
dụng các mô hình đã biết để giải nhiều dạng bài, mà còn yêu cầu khi gặp bài
mới vượt ra ngoài phạm vi của các mô hình kinh điển, ta phải biết phân tích
tình huống mới một cách sáng tạo và giải quyết vấn đề mới. Vì vậy đây là một
loại bài khá thách thức.

Chẳng hạn bài ``Mạng điện trở'' trong kỳ tuyển chọn đội Trung Quốc dự IOI
1998 là một ví dụ điển hình; đề bài được đặt trong phụ lục của bài gốc.

Đây là một bài lấy nguồn từ điện học vật lý. Khi xuất hiện trong thi tin học,
từ đầu đến cuối nó đều mới, không có bất kỳ mô hình kinh nghiệm nào sẵn để
tham khảo. Nếu bài này có một thuật toán đặc thù nào, thì đó là dùng vài công
thức vật lý cơ bản của mạch nối tiếp và song song đơn giản, từng bước gộp cặp
những điện trở có thể gộp, cho đến khi cuối cùng chỉ còn duy nhất một điện
trở, tức tổng điện trở của mạch ngoài. Đây là ``suy tiến''. Sau đó ta lần
ngược theo các bước ``suy tiến'', cho đến khi tổng điện trở ấy được mở rộng
trở lại thành mạng điện trở ban đầu; trong quá trình truy ngược, ta tính được
dòng điện và điện áp trên mọi điện trở. Đây là ``suy ngược''.

Nhưng để thực hiện ``hai khúc'' ấy không phải việc dễ. Trở ngại lớn nhất là
không có sẵn cấu trúc dữ liệu chuyên dụng để mô tả mạch điện. Vì vậy, điểm
then chốt của bài nằm ở việc làm thế nào dùng ngôn ngữ lập trình hữu hạn để
lưu trữ, biểu diễn và thao tác trên mạng điện trở đan xen ngang dọc này, sao
cho vừa tiện lợi, nhanh chóng, vừa đảm bảo không sinh mơ hồ. Việc thiết kế
một bộ mô hình cấu trúc dữ liệu đẹp cũng mở ra cho ta một không gian rộng lớn
để phát huy năng lực sáng tạo và tưởng tượng. Có lẽ việc kiểm tra mức độ linh
hoạt trong vận dụng tri thức và tư duy sáng tạo của thí sinh cũng chính là ý
đồ của người ra bài.

\subsubsection{Chiến lược: bốc thuốc đúng bệnh}

Vì những mô hình được đưa vào bởi mở hóa biến đổi muôn hình vạn trạng và
không có mẫu giải cố định nào để lần theo, chỉ có ứng biến theo tình huống và
bốc thuốc đúng bệnh mới là phương pháp chung duy nhất để giải loại bài này.
Bài gốc kèm chương trình tham khảo \texttt{resistor.pas}.

\subsubsection{Tiểu kết}

Mô hình mở hóa liên tục xuất hiện trong thi đấu là sản phẩm tất yếu của việc
tin học ngày càng gắn với thực tiễn. Sự phong phú của thế giới khách quan
quyết định rằng, khi dùng lập trình để giải quyết vấn đề thực tế, những vấn
đề ta gặp sẽ hết sức đa dạng. Vì vậy, chưa từng có một mô hình ``vạn năng''
nào có thể thử đâu thắng đó. Ta chỉ có thể xuất phát từ tính đặc thù của vấn
đề cụ thể, phân tích cụ thể từng tình huống cụ thể, mới có thể tìm được cách
giải tối ưu độc nhất của bài toán đó.

\section{Tổng kết}

Ẩn hóa, đa chiều hóa và mở hóa là ba biểu hiện lớn của tính linh hoạt trong
mô hình hóa toán học ở thi tin học. Cần thấy rằng ba xu hướng này không tồn
tại độc lập với nhau, mà hòa quyện và xuyên suốt các bài thi. Chính bài
``Tuyến du lịch tối ưu'' đã nhắc trong phần ẩn hóa là một ví dụ. Trong quá
trình ẩn mô hình, người ra bài thực ra đã dùng hiện tượng hai chiều để mô tả
một mô hình một chiều. Điều này khác với đa chiều hóa đã nói ở trên, vì
``chiều'' ở đây thật ra không phải là chiều thực, nhưng ta vẫn nên thừa nhận
rằng trong đó cũng gửi gắm tư tưởng đa chiều. Vì vậy, ta nên nhìn ba xu
hướng này bằng con mắt liên hệ, không tách rời. Thường chính sự kết hợp hài
hòa của cả ba mới đem lại cho chúng ta những bài thi mới hơn, hay hơn và giàu
tính sáng tạo hơn.

Còn với tư cách thí sinh, tính linh hoạt của mô hình hóa toán học trong thi
tin học chắc chắn đặt ra yêu cầu cao hơn đối với phẩm chất toàn diện của
chúng ta. Tất nhiên, ta nên ý thức rằng đây không chỉ là yêu cầu của kỳ thi,
mà còn là kỳ vọng tha thiết và lời kêu gọi nóng hổi của thời đại đối với thế
hệ thanh thiếu niên chúng ta. Tham gia thi tin học, chủ động bồi dưỡng các
năng lực và phẩm chất của mình trên nhiều mặt, đặc biệt là năng lực mô hình
hóa toán học và năng lực lập trình, cũng chính là hành động thực tế quan
trọng để chúng ta không ngừng theo đuổi sự phát triển và hoàn thiện bản thân,
đồng thời nỗ lực chuẩn bị cho tương lai.

\section*{Tài liệu tham khảo}

\begin{enumerate}
  \item Phân tích đề thi tin học (máy tính) thanh thiếu niên quốc tế và trong
  nước, 1994--1995.
  \item Phân tích đề thi tin học (máy tính) thanh thiếu niên quốc tế và trong
  nước, 1992--1993.
\end{enumerate}

\end{document}
